\chapter{Arthroscopy of the Knee}
\index{Arthroscopy of the Knee}

\section*{Definition and Overview}
Arthroscopy of the knee is a keyhole operation that allows a surgeon to look inside the knee joint with a small camera and treat a wide range of injuries and conditions through incisions of only a few millimetres. It is one of the most commonly performed orthopaedic procedures in the world. The technique is used to repair or trim a torn meniscus, reconstruct a torn anterior cruciate ligament (ACL), remove loose fragments, smooth damaged cartilage, or wash out an infected joint. Because the incisions are small, there is less tissue damage, less pain, and a faster recovery than with traditional open surgery.

\section*{Epidemiology}
Knee arthroscopy is performed across all age groups. Meniscal and ligament tears are common in younger, active people who injure the knee in sport, while older adults more often undergo arthroscopy for degenerate or damaged cartilage. ACL injuries are among the most frequent serious knee injuries in athletes, particularly in sports involving pivoting and cutting movements such as football, basketball, and netball. The overall incidence of knee injuries rises with sporting participation and with age, making knee arthroscopy a major component of orthopaedic workload.

\section*{Aetiology and Pathophysiology}
\begin{itemize}
    \item \textbf{Meniscal tears:} twisting injuries while the foot is planted, or age-related degeneration that weakens the meniscus so that it tears with minor stress
    \item \textbf{ACL tears:} sudden deceleration, pivoting, or landing awkwardly, often accompanied by a popping sensation and immediate swelling
    \item \textbf{Other ligament injuries:} sprains of the medial and lateral collateral ligaments, which often settle without surgery
    \item \textbf{Articular cartilage damage:} injury, wear, or the development of loose flaps of cartilage
    \item \textbf{Loose bodies:} free fragments of bone or cartilage within the joint that cause catching and locking
    \item \textbf{Infection or inflammation:} bacterial or crystal-induced inflammation of the joint lining requiring washout
    \item \textbf{Plica syndrome and other synovial problems:} thickened folds of joint lining that catch within the joint
\end{itemize}

\section*{Clinical Features}
\begin{itemize}
    \item Knee pain, swelling, or stiffness after injury or with activity
    \item Clicking, catching, or locking of the knee during movement
    \item A feeling of instability, or the knee giving way under load
    \item Difficulty fully straightening or bending the knee
    \item Reduced ability to weight-bear or to participate in sport
    \item After surgery: mild pain, swelling, and bruising around the incisions for a few days
\end{itemize}

\section*{Differential Diagnosis}
Conditions that can cause similar knee symptoms include:
\begin{itemize}
    \item Osteoarthritis and rheumatoid arthritis
    \item Patellofemoral pain syndrome, a common cause of anterior knee pain
    \item Bursitis or tendinopathy around the knee
    \item Ligament sprains that heal with non-operative treatment
    \item Referred pain from the hip or lumbar spine
    \item Stress fracture or other bone pathology
    \item Crystal arthropathies such as gout and pseudogout
\end{itemize}

\section*{When to Seek Medical Care}
See a doctor promptly after a knee injury if the knee becomes very swollen, cannot bear weight, locks, or gives way, or if pain persists after a few days of rest.

\textbf{Contact your local emergency services (or attend an emergency department) for:}
\begin{itemize}
    \item A knee that is obviously deformed or dislocated
    \item An injury followed by numbness, coldness, or loss of pulse in the leg
    \item Sudden inability to move the knee or bear any weight
    \item After arthroscopy: severe pain, fever, spreading redness, or calf swelling and pain that may indicate a blood clot
\end{itemize}

\section*{Diagnosis and Investigations}
\begin{itemize}
    \item \textbf{History:} how the injury occurred and what the knee does during daily activity and sport
    \item \textbf{Examination:} assessment of swelling, tenderness, range of movement, and stability using tests such as the Lachman and McMurray tests
    \item \textbf{X-ray:} to exclude fractures and assess alignment and the degree of arthritis
    \item \textbf{MRI:} the most sensitive imaging for meniscal, ligament, and cartilage injuries
    \item \textbf{Ultrasound:} occasionally used to assess tendons or fluid collections
    \item \textbf{Arthroscopy:} direct visual inspection, which both confirms the diagnosis and allows treatment in the same procedure
\end{itemize}

\section*{Management}
\begin{itemize}
    \item \textbf{Non-operative care:} rest, ice, compression, elevation, analgesia, and physiotherapy for many ligament sprains and minor meniscal injuries
    \item \textbf{Partial meniscectomy:} trimming of the torn portion of a meniscus that is unlikely to heal
    \item \textbf{Meniscal repair:} stitching of a tear in the well-vascularised outer zone of the meniscus, which has a good healing potential
    \item \textbf{ACL reconstruction:} rebuilding the ligament with a tendon graft, usually followed by a lengthy rehabilitation programme
    \item \textbf{Debridement and removal of loose bodies:} for damaged cartilage and free fragments within the joint
    \item \textbf{Washout:} arthroscopic irrigation of an infected knee together with antibiotics
    \item \textbf{Rehabilitation:} structured physiotherapy, which is essential for regaining strength, stability, and full range of movement
\end{itemize}

\section*{Complications}
\begin{itemize}
    \item Infection in the knee joint or at the incisions
    \item Bleeding or swelling within the joint
    \item Deep vein thrombosis or pulmonary embolism after surgery
    \item Stiffness, particularly after ligament reconstruction
    \item Ongoing pain, locking, or instability if the injury was extensive
    \item Damage to nerves, blood vessels, or other joint structures during surgery
    \item An increased risk of earlier arthritis after meniscal removal or significant cartilage injury
\end{itemize}

\section*{Prognosis and Outcomes}
Most people recover well, and many simple procedures allow a return to full activity within weeks. Meniscal repairs and ACL reconstructions in appropriately selected patients generally have good outcomes, although full return to sport often takes nine to twelve months after ligament reconstruction. Recovery depends on the extent of the injury and on committed rehabilitation. In many cases the results are excellent, but some residual discomfort or a slightly earlier onset of arthritis may persist after more extensive damage.

\section*{Prevention and Self-Care}
\begin{itemize}
    \item Strengthen the thigh and hip muscles to protect the knee during sport and daily activity
    \item Warm up properly and use correct technique, especially in pivoting sports
    \item Choose footwear and playing surfaces appropriate to the activity
    \item Maintain a healthy weight to reduce stress on the knees
    \item Follow the prescribed rehabilitation exercises after surgery
    \item Return to sport only once strength, range, and stability have been regained
\end{itemize}

\section*{Sources and Further Reading}
The following organisations provide reliable information on knee arthroscopy for patients and healthcare students.

\begin{itemize}
    \item NHS. \textit{Arthroscopy of the knee: overview and recovery.} https://www.nhs.uk/conditions/arthroscopy/
    \item American Academy of Orthopaedic Surgeons (AAOS). \textit{Knee arthroscopy.} https://orthoinfo.aaos.org/en/treatment/knee-arthroscopy/
    \item Mayo Clinic. \textit{Arthroscopy overview.} https://www.mayoclinic.org/tests-procedures/arthroscopy/
    \item MedlinePlus. \textit{Knee arthroscopy.} https://medlineplus.gov/ency/article/007207.htm
    \item Versus Arthritis. \textit{Knee pain and joint surgery information.} https://www.versusarthritis.org/about-arthritis/conditions/knee-pain/
    \item MSD Manual. \textit{Knee pain and arthroscopy.} https://www.msdmanuals.com/professional/musculoskeletal-and-connective-tissue-disorders/approach-to-the-patient-with-joint-disease/knee-pain
\end{itemize}
